% main.tex 主编译文件 编译引擎：XeLaTeX
\documentclass[12pt,a4paper,heading]{ctexart}
% 屏蔽xeCJK字体重定义冗余警告
\xeCJKsetup{silent=redefinefont}

% ====================== 字体设置 ======================
\usepackage{fontspec}
\setmainfont{Times New Roman}
% 弱化伪加粗，减少黑体缺失警告
\setCJKmainfont{SimSun}[AutoFakeBold=0.6]

% ====================== 全局配色 & 表格支持 ======================
\usepackage[table]{xcolor}
\usepackage{multirow}
\definecolor{subtitleblue}{RGB}{0,60,130}
\definecolor{keyred}{RGB}{200,0,0}

% ====================== 页面基础版式 ======================
\usepackage{geometry}
\geometry{left=2.4cm,right=2.4cm,top=2.6cm,bottom=2.6cm,includeheadfoot}

\usepackage{setspace}
\onehalfspacing
\usepackage{indentfirst}
\setlength{\parindent}{2em}

% ====================== CTEX官方标题配置（安全不破坏目录） ======================
\ctexset{
    section = {
        format = \centering\heiti\zihao{3},
        beforeskip = 20pt,
        afterskip = 14pt,
    },
    subsection = {
        format = \heiti\zihao{4}\textcolor{subtitleblue},
        beforeskip = 14pt,
        afterskip = 10pt,
    },
    subsubsection = {
        format = \heiti\zihao{4},
        beforeskip = 10pt,
        afterskip = 8pt,
    }
}

% ====================== 目录配置【恢复原生点线、紧凑靠左收拢，标准tocloft语法】 ======================
\usepackage{tocloft}
\renewcommand{\contentsname}{\Large\heiti 目 录}
% 1. 删除清空点线的代码，恢复默认圆点连接线
% 2. 目录字体设置
\renewcommand{\cftsecfont}{\heiti}
\renewcommand{\cftsecpagefont}{\heiti}
\renewcommand{\cftsubsecfont}{}
\renewcommand{\cftsubsubsecfont}{}

% 3. 缩进整体收窄，层级横向距离变小（紧凑效果保留）
\cftsetindents{section}{0em}{1.2em}
\cftsetindents{subsection}{1.2em}{1.2em}
\cftsetindents{subsubsection}{2.4em}{1.2em}

% 4. 各级编号占用宽度缩小，减少空白
\setlength{\cftsecnumwidth}{1.2em}
\setlength{\cftsubsecnumwidth}{1.4em}
\setlength{\cftsubsubsecnumwidth}{1.6em}

% 5. 收紧右侧留白、缩小页码预留宽度（官方正确语法）
\cftsetpnumwidth{1.2cm}
\cftsetrmarg{1.4cm}

% ====================== 页眉页脚 ======================
\usepackage{fancyhdr}
\setlength{\headheight}{16pt}
\pagestyle{fancy}
\fancyhf{}
\fancyhead[L]{\small 毛概课程笔记}
\fancyhead[R]{\small 浮萍}
\fancyfoot[C]{\thepage}
\renewcommand{\headrulewidth}{0.4pt}
\renewcommand{\footrulewidth}{0.3pt}
\fancypagestyle{empty}{\fancyhf{}\renewcommand{\headrulewidth}{0pt}\renewcommand{\footrulewidth}{0pt}}

% ====================== 表格、列表工具 ======================
\usepackage{array}
\usepackage{tabularx}
\usepackage{enumitem}
\usepackage{pifont}

\setlist[itemize]{leftmargin=2em,itemsep=0.4em}
\setlist[enumerate]{leftmargin=2em}

\renewcommand{\arraystretch}{1.6}
\setlength{\tabcolsep}{6pt}
\newlist{dotlist}{enumerate}{1}
\setlist[dotlist]{label=\ding{182},leftmargin=2em,itemsep=0.5em}
% 表格居中环境
\newenvironment{notetable}{\begin{center}\renewcommand{\arraystretch}{1.6}}{\end{center}}

% ====================== 超链接、书签 ======================
\usepackage{hyperref}
\hypersetup{
    colorlinks=true,
    linkcolor=black,
    urlcolor=blue,
    citecolor=black,
    pdfborder={0 0 0},
    pdfauthor={浮萍},
    pdftitle={毛概课程完整笔记}
}
\usepackage{bookmark}

% ====================== 自定义快捷命令 ======================
% 取消段落首缩进
\newcommand{\noparindent}{\setlength{\parindent}{2em}}
% 红色黑体重点标注
\newcommand{\keytext}[1]{\textcolor{keyred}{\heiti #1}}

\begin{document}

% 封面无页眉页脚
\thispagestyle{empty}
% 2.cover.tex 《毛泽东思想和中国特色社会主义理论体系概论》课程笔记封面
\begin{titlepage}
    \thispagestyle{empty}
    \centering

    % 顶部留白
    \vspace*{1.8cm}

    % 课程名称：使用黑体加粗，规避SimSun粗体缺失警告
    {\Large \heiti 毛泽东思想和\\[0.3em]
    中国特色社会主义理论体系概论\par}

    \vspace{0.6cm}
    \rule{0.68\textwidth}{0.6pt}
    \vspace{0.7cm}

    % 主标题
    {\LARGE \heiti 课程笔记\par}

    \vspace{0.6cm}

    % 副标题
    {\large 新时代青年传承红色文化的路径探析\par}

    \vspace{1.4cm}

    % 引言：取消宋体斜体，改用常规宋体，消除字体警告
    \begin{minipage}{0.75\textwidth}
        \centering
        {\small
        “红色文化是中国共产党百年奋斗历程中积淀而成的精神财富，\\
        也是新时代青年坚定理想信念、涵养家国情怀的重要源泉。”\par}
    \end{minipage}

    \vfill

    % 下方信息分隔线
    \rule{0.48\textwidth}{0.5pt}
    \vspace{0.9cm}

    {\normalsize
    \renewcommand{\arraystretch}{1.35}
    \begin{tabular}{rl}
        作者：   & 浮萍 \\
        版本：   & 1.0 完整版 \\
        邮箱：   & \texttt{3071974740@qq.com} \\
        GitHub： & \texttt{https://github.com/Duckweed-yhb} \\
        日期：   & \today \\
    \end{tabular}
    \par}

    \vspace{1.2cm}
\end{titlepage}

\clearpage
\setcounter{page}{1}
\pagestyle{fancy}

% 生成目录
\tableofcontents
\clearpage

% 依次导入章节
% 子文件 3.导论，仅input导入，无document
\noparindent
\section{导论 马克思主义中国化时代化的历史进程和理论成果}
\vspace{1em}

\subsection{两个相结合}
把马克思主义基本原理同中国具体实际相结合、同中华优秀传统文化相结合。

\vspace{1em}
\subsection{马克思主义中国化时代化的提出及其原因和内涵}
\begin{enumerate}
    \item \keytext{提出}：1938年，毛泽东在党的六届六中全会上作的题为《论新阶段》的政治报告中最先提出了“马克思主义中国化”这个命题。这是党首次明确提出“马克思主义中国化”的重大命题。马克思主义中国化同时包含着马克思主义时代化的意蕴。

    \item \keytext{原因}：推进马克思主义中国化时代化，是马克思主义理论本身发展的内在要求；是解决中国实际问题的客观需要。

    \item \keytext{科学内涵}：毛泽东所说的“解决中国问题”和“创造些新的东西”，实际上指向马克思主义中国化时代化的科学内涵，具体包含三层意思：
    \begin{itemize}
        \item 运用马克思主义的立场、观点和方法，观察时代、把握时代、引领时代，解决中国革命、建设、改革中的实际问题。
        \item 总结和提炼中国革命、建设、改革的实践经验，并将其上升为理论，不断丰富和发展马克思主义的理论宝库，赋予马克思主义以新的时代内涵。
        \item 运用中国人民喜闻乐见的民族语言来阐述马克思主义，使其植根于中华优秀传统文化的土壤之中，具有中国特色、中国风格、中国气派。
    \end{itemize}
\end{enumerate}

\vspace{1em}
\subsection{准确把握马克思主义中国化时代化的要求}
\begin{itemize}
    \item 推进马克思主义中国化时代化是一个追求真理、揭示真理、笃行真理的过程。
    \item 准确把握马克思主义中国化时代化，要做到坚持马克思主义与发展马克思主义相统一。
    \item 坚持和发展马克思主义，必须同中国具体实际相结合、同中华优秀传统文化相结合。
\end{itemize}

\vspace{1em}
\subsection{马克思主义中国化时代化的四个历史进程}
\begin{notetable}
\textbf{\large 马克思主义中国化时代化四个历史进程}
\par\vspace{0.5em}
\begin{tabular}{|>{\centering\arraybackslash}m{1.6cm}|>{\centering\arraybackslash}m{3.6cm}|m{6.2cm}|}
\hline
时间 & 历史分期 & 历史成就 \\ \hline
1921-1949 年 & 新民主主义革命时期 & 创造了新民主主义革命的伟大成就，为实现中华民族伟大复兴创造根本社会条件 \\ \hline
1949-1978 年 & 社会主义革命和建设时期 & 创造了社会主义革命和建设的伟大成就，为实现中华民族伟大复兴奠定根本政治前提和制度基础 \\ \hline
1978-2012 年 & 改革开放和社会主义现代化建设新时期 & 创造了改革开放和社会主义现代化建设的伟大成就，为实现中华民族伟大复兴提供充满活力的体制保证与物质条件 \\ \hline
2012 年至今 & 中国特色社会主义新时代 & 创造新时代伟大成就，提供更完善制度、更坚实物质、更主动精神力量 \\ \hline
\end{tabular}
\end{notetable}

\vspace{1em}
\subsection{马克思主义中国化时代化的理论成果及其相互关系}
\begin{enumerate}
    \item \keytext{主要理论成果}：在百年奋斗历程中，党先后形成毛泽东思想、邓小平理论、“三个代表”重要思想、科学发展观、习近平新时代中国特色社会主义思想，一脉相承、与时俱进。

    \item \keytext{成果内在关系}：全部理论成果整体呈现\keytext{一脉相承又与时俱进}特征：
    \begin{itemize}
        \item 毛泽东思想蕴含立场观点方法，为后续理论筑牢根基；
        \item 中特理论体系立足新实践，丰富发展毛泽东思想；
        \item 全部成果均是马列主义在中国的运用发展，是长期指导思想、共同思想基础。
    \end{itemize}
\end{enumerate}
% 子文件：4.毛泽东思想及其历史地位，仅input导入
\noparindent
\section{毛泽东思想及其历史地位}

\subsection{毛泽东思想的形成和发展}
\subsubsection{毛泽东思想形成发展的历史条件}
\begin{enumerate}
    \item \textbf{时代背景}：19世纪末20世纪初，世界进入帝国主义和无产阶级革命时代，战争与革命成为时代主题。
    \item \textbf{实践基础}：中国共产党领导人民进行的革命和建设的实践。
\end{enumerate}

\subsubsection{毛泽东思想形成发展的过程}
\paragraph{教材梳理版}
\textbf{时代主题：战争与革命}
\begin{enumerate}
    \item \textbf{开始萌芽：大革命时期}
    \begin{itemize}
        \item 标志性著作：《中国社会各阶级的分析》《湖南农民运动考察报告》等——分析社会各阶级在革命中的地位和作用。
        \item 核心思想：明确无产阶级领导农民是革命核心问题，新民主主义革命思想萌芽。
    \end{itemize}

    \item \textbf{初步形成：土地革命战争时期}
    \begin{itemize}
        \item 标志性著作：《中国的红色政权为什么能够存在?》《井冈山的斗争》《星星之火，可以燎原》《反对本本主义》。
        \item 核心思想：提出农村包围城市、武装夺取政权道路，标志理论初步形成。
    \end{itemize}
    \vspace{0.4em}
    \noindent\sloppy\colorbox{orange!10}{
        \begin{minipage}{0.92\textwidth}
            \textbf{点拨：}中国革命道路：农村建立根据地积蓄力量，时机成熟夺取全国政权。
        \end{minipage}
    }
    \vspace{0.4em}

    \item \textbf{\color{red}成熟：遵义会议后至全民族抗战时期}
    \begin{itemize}
        \item 理论著作：《实践论》《矛盾论》《中国革命和中国共产党》《新民主主义论》等；
        \item 关键节点：1945党的七大将毛泽东思想写入党章，确立指导地位。
    \end{itemize}

    \item \textbf{\color{red}继续发展：解放战争与建国后}
    \begin{itemize}
        \item 标志性著作：七届二中全会报告、《论人民民主专政》《论十大关系》《关于正确处理人民内部矛盾的问题》；
        \item 核心任务：提出马列主义与中国实际“第二次结合”，探索社会主义建设道路。
    \end{itemize}
\end{enumerate}

\paragraph{两次结合拓展}
\begin{itemize}
    \item \textbf{第一次结合}：新民主主义革命阶段，走出革命道路；
    \item \textbf{第二次结合}：社会主义建设阶段，探索建设道路。
\end{itemize}
\vspace{0.3em}
\textbf{“第二次结合”详细阐述}
建国后党完成社会主义改造，进入大规模建设阶段。毛泽东提出把马列主义同中国实际第二次结合，独立探索适合国情的社会主义建设道路，形成人民民主专政、两类矛盾、工业化等理论成果。

\subsection{毛泽东思想的主要内容和活的灵魂}
\subsubsection{毛泽东思想的主题和主要内容}
\begin{enumerate}
    \item \textbf{主题}：中国革命和建设。
    \item \textbf{六大核心内容}：
    \begin{enumerate}
        \item 新民主主义革命理论
        \item 社会主义革命和建设理论
        \item 革命军队与军事战略理论
        \item 政策和策略理论
        \item 思想政治与文化工作理论
        \item 党的建设理论
    \end{enumerate}
\end{enumerate}
\vspace{0.3em}
补充：\keytext{统一战线、武装斗争、党的建设}是革命三大法宝。

\paragraph{拓展：窑洞对}
\textbf{跳出历史周期率两条路径}
\begin{enumerate}
    \item 人民监督政府
    \item 党的自我革命
\end{enumerate}

\subsubsection{毛泽东思想活的灵魂}
\begin{enumerate}
    \item \textbf{总述}：实事求是、群众路线、独立自主；实事求是为精髓。
\end{enumerate}
\vspace{0.4em}
\noindent\sloppy\colorbox{orange!10}{
    \begin{minipage}{0.92\textwidth}
        \textbf{点拨：}1981年历史决议首次概括毛泽东思想活的灵魂。
    \end{minipage}
}
\vspace{0.4em}

\begin{enumerate}[resume]
    \item \textbf{\color{red}实事求是}
    \begin{itemize}
        \item 定位：思想路线、理论精髓；
        \item 内涵：一切从实际出发，理论联系实际，实践检验真理；
        \item 出处：《实践论》《矛盾论》。
    \end{itemize}

    \item \textbf{\color{red}群众路线}
    \begin{itemize}
        \item 定位：党的生命线、根本工作路线；
        \item 内涵：一切为群众、一切依靠群众，从群众中来、到群众中去。
    \end{itemize}

    \item \textbf{\color{red}独立自主}
    \begin{itemize}
        \item 定位：立党立国根本原则；
        \item 出处：1930《反对本本主义》；
        \item 内涵：立足本国力量，自主发展，合理借鉴外部经验。
    \end{itemize}
\end{enumerate}
\vspace{0.4em}
\noindent\sloppy\colorbox{orange!10}{
    \begin{minipage}{0.92\textwidth}
        \textbf{点拨：}《反对本本主义》：中国革命斗争的胜利要靠中国同志了解中国情况。
    \end{minipage}
}
\vspace{0.4em}

\subsection{毛泽东思想的历史地位}
\subsubsection{定义}
毛泽东思想是马列主义在中国的运用发展，是中国革命建设正确理论总结，是党的集体智慧结晶。

\subsubsection{历史地位}
\begin{enumerate}
    \item 马克思主义中国化第一次历史性飞跃理论成果；
    \item 中国革命和建设的科学指南；
    \item 党和人民永久宝贵精神财富。
\end{enumerate}

\subsubsection{正确认识毛泽东与毛泽东思想}
\begin{enumerate}
    \item 评价历史人物必须放在所处时代历史条件下；
    \item 领袖是人不是神，既不能神化崇拜，也不能全盘否定陷入历史虚无；
    \item 总体评价：功绩第一位，错误第二位。
\end{enumerate}

\subsection{本章小结}
毛泽东思想是马克思主义中国化首个重大理论成果，形成于土地革命、成熟于抗战、建国后持续发展。六大理论体系为核心，实事求是、群众路线、独立自主是活的灵魂，长期作为党的指导思想。
% 子文件：5.新民主主义革命理论，仅input导入
\noparindent
\section{新民主主义革命理论}

\subsection{近代中国国情}
\begin{enumerate}
    \item \textbf{认清国情}：认识解决中国革命一切基本依据；
    \item \textbf{基本国情}：半殖民地半封建社会；
    \item \textbf{主要矛盾}：帝国主义与中华民族（最主要）、封建主义与人民大众；
    \item \textbf{根本任务}：推翻三座大山，扫清复兴发展障碍；
    \item \textbf{革命性质}：资产阶级民主革命，非无产阶级社会主义革命。
\end{enumerate}

\subsection{近代中国革命时代特征}
\begin{enumerate}
    \item 十月革命后，中国民主革命转为世界无产阶级革命一部分；
    \item 五四运动标志进入新民主主义革命阶段。
\end{enumerate}

\subsection{新民主主义革命总路线}
\begin{enumerate}
    \item 1939《中国革命和中国共产党》首次提出；
    \item 1948《在晋绥干部会议上的讲话》完整概括：无产阶级领导，人民大众，反帝反封建反官僚资本主义；
    \item 作用：明确革命对象、动力、领导权、性质、前途。
\end{enumerate}

\subsection{新民主主义革命的动力}
\paragraph{人民大众组成}工人阶级（领导）、农民阶级（主力军）、城市小资产阶级、民族资产阶级。

\subsection{新民主主义革命领导力量}
\begin{enumerate}
    \item 无产阶级领导权是新旧民主革命根本区分标志；
    \item 无产阶级自身特点决定其领导地位，是革命胜利根本保证。
\end{enumerate}
\vspace{0.4em}
\noindent\sloppy\colorbox{orange!10}{
    \begin{minipage}{0.92\textwidth}
        \textbf{点拨：}实现领导权五大条件：\ding{172}联合斗争取胜；\ding{173}工农联盟为关键；\ding{174}又联合又斗争；\ding{175}革命武装作支柱；\ding{176}加强党建为根本。
    \end{minipage}
}
\vspace{0.4em}

\paragraph{核心概念对照表}
\begin{center}
\begin{tabularx}{0.95\textwidth}{|>{\centering\arraybackslash}X|>{\centering\arraybackslash}X|}
\hline
\rowcolor{orange!10} 关键命题 & 核心内容 \\
\hline
革命首要问题 & 分清敌友 \\
\hline 革命中心/理论核心 & 无产阶级领导权 \\
\hline 革命基本问题 & 农民问题 \\
\hline 革命主要内容 & 没收地主土地归农民 \\
\hline 革命实质 & 中共领导下农民革命 \\
\hline
\end{tabularx}
\end{center}

\subsection{新民主主义革命性质与前途}
\begin{enumerate}
    \item 属于新式资产阶级民主革命，区别旧民主革命；
    \item 四大新特征：世界无产革命一部分、中共领导、马列指导、前途社会主义。
\end{enumerate}

\begin{enumerate}[resume]
    \item 新民主主义革命与社会主义革命的区别——二者性质不同。
    \item 新民主主义革命与社会主义革命的联系——紧密衔接，民主革命内含社会主义因素。
    \begin{itemize}
        \item 论断：民主主义革命是社会主义革命的必要准备，社会主义革命是民主主义革命的必然趋势。
        \item 中国革命必须分“两步走”：第一步民主主义革命，第二步社会主义革命。
    \end{itemize}
    \item 认清二者关系的重要性：只有分清区别、看清联系，才能正确领导中国革命。
\end{enumerate}

\vspace{0.4em}
\paragraph{新旧民主革命与社会主义革命区分表}
\begin{center}
\begin{tabularx}{0.97\textwidth}{|>{\centering\arraybackslash}X|>{\centering\arraybackslash}X|>{\centering\arraybackslash}X|>{\centering\arraybackslash}X|>{\centering\arraybackslash}X|}
\hline
\rowcolor{orange!10} 类别 & 革命性质 & 资本政策 & 国家前途 & 二者关系 \\
\hline
新民主主义革命 & 资产阶级民主革命 & 没收官僚、保护民族资本 & 各革命阶级联合专政 & \multirow{2}{=}两步走：民主革命是准备，社会主义是必然趋势 \\
\cline{1-4}
社会主义革命 & 无产阶级革命 & 消灭私有制 & 无产阶级专政 & \\
\hline
\end{tabularx}
\end{center}

\vspace{0.4em}
\noindent\sloppy\colorbox{orange!10}{
    \begin{minipage}{0.92\textwidth}
        \textbf{点拨：}
        \ding{172} “一次革命论”（左倾无间断革命论）：只看见联系、混淆阶段区别，反帝反封建同时反对民族资产阶级，否定革命阶段性；
        \ding{173} “二次革命论”（右倾错误）：只看见区别、割裂阶段联系，放弃无产阶级对民主革命的领导权。
        \par 切勿混淆左倾“无间断革命论”与马克思主义“不断革命论”。
    \end{minipage}
}
\vspace{0.4em}

\subsection{新民主主义革命的基本纲领}
新民主主义基本纲领是总路线的具体展开，为革命指明具体奋斗目标。
\begin{enumerate}
    \item \textbf{政治纲领}：推翻帝国主义、封建主义统治，建立无产阶级领导、工农联盟为基础、各革命阶级联合专政的新民主主义共和国；属于过渡性国家形态。
    \item \textbf{经济纲领}三大内容：
    \begin{itemize}
        \item 没收封建地主土地归农民所有（新民主主义革命核心内容）；
        \item 没收官僚垄断资本归新民主主义国家所有（具备双重性质）；
        \item 保护民族工商业（由落后国情与革命性质决定，独有特色）。
    \end{itemize}
    三大经济纲领是新民主主义共和国经济建设根本方针。
\end{enumerate}
\vspace{0.4em}
\noindent\sloppy\colorbox{orange!10}{
    \begin{minipage}{0.92\textwidth}
        \textbf{点拨：}反对官僚资本主义，不是因为它属于资本主义，而是官僚资本依附帝国主义、封建地主势力，是买办封建式国家垄断资本，是国民党政权经济根基，严重阻碍中国发展。
    \end{minipage}
}
\vspace{0.4em}

\begin{enumerate}[resume]
    \item \textbf{文化纲领}：无产阶级领导、人民大众、反帝反封建的文化，即民族的、科学的、大众的文化。
    \begin{itemize}
        \item 民族的文化：反帝独立，保有民族风格特色，兼容外来进步文化；
        \item 科学的文化：反对封建迷信，坚持实事求是、理论实践统一；
        \item 大众的文化：服务工农劳动群众，面向大众、提升群众，具有民主属性。
    \end{itemize}
\end{enumerate}

\subsection{新民主主义革命道路：农村包围城市、武装夺取政权}
\subsubsection{革命道路的提出过程}
\begin{enumerate}
    \item \textbf{初步形成}：1927大革命失败后，毛泽东创建井冈山根据地，撰写《中国的红色政权为什么能够存在?》《井冈山的斗争》《星星之火，可以燎原》，提出“工农武装割据”，初步构建农村包围城市理论。
    \item \textbf{正式确立}：党的六届六中全会明确：党的任务不是长期合法斗争、先占城市后取乡村，而是走农村包围城市的相反道路。
\end{enumerate}

\subsubsection{道路形成的历史必然性}
由近代中国半殖民地半封建特殊国情决定：
\begin{enumerate}
    \item 国内无民主制度、受封建压迫；外部无民族独立、受帝国主义压迫，革命唯一斗争形式为武装斗争，核心组织是人民军队。
    \item 近代中国农业大国，农民占人口绝大多数，是无产阶级可靠同盟军、革命主力军；党必须深入农村解决土地问题，依靠农民建立革命根基，最终夺取全国政权。
\end{enumerate}

\end{document}